% !TeX root = ../教材题库_全部条目_转写工作稿.tex
% Source: XJ-B1-C06-HW-P254-Y0953
% 第6章 / 6.4 用样本估计总体 / 教材 p.254 / PDF 第261页
% OCR未校对：公式与文字均可能有误；未完成数学TeX转写；禁止直接发给学生
% 原题及相邻原解截图（核对用）：../题目原图/XJ-B1-C06-HW-P254-Y0953-01.webp
% 本文件只在转写完成后作为题干使用；OCR 条目尚不是可靠题干。
\begin{bankitem}{习题6.4 第2题}{第6章 / 6.4 用样本估计总体 / 教材 p.254 / PDF 第261页}{XJ-B1-C06-HW-P254-Y0953}{OCR未校对：公式与文字均可能有误；未完成数学TeX转写；禁止直接发给学生}
\begin{ocrdraft}
[教材p.254 / PDF第261页]
\par
2．下面是国外某城市某日在不同观测点对细颗粒物（PM25）浓度的观测值：
\par
275
\par
268
\par
237
\par
208
\par
225
\par
396
\par
168
\par
199
\par
巧7
\par
166176173188
\par
221
\par
176
\par
巧9
\par
168
\par
巧0
\par
173
\par
198
\par
177
\par
129
\par
14163141142
\par
巧7
\par
142
\par
112
\par
136
\par
140
\par
166
\par
102
\par
110
\par
98
\par
（单位：gg/m）
\par
（1）数据中有无众数？
\par
（2）计算数据的中位数与均值，它们相等吗？
\par
（3）若数据中的最大值比现有的最大值多25，数据的极差、中位数、众数、平
\par
均数发生改变了吗？
\par
（4）根据上述数据，估计该市当日细颗粒物浓度的中位数、众数．
\end{ocrdraft}
\end{bankitem}
